\BourbakiExercises{习题}

\begin{enumerate}
\def\labelenumi{\arabic{enumi})}
\tightlist
\item
  设 p 是一个素数。用 C 表示形如 \((\mathbf{Z}/p^{2}\mathbf{Z})^{n} \oplus (\mathbf{Z}/p^{3}\mathbf{Z})^{m}\)（\(n, m \in N\)）的 Z-模的类所成之集。
\end{enumerate}

\begin{enumerate}
\def\labelenumi{\alph{enumi})}
\item
  证明 \(\mathcal{C}\) 是加性的但不是遗传的。最小的遗传集 \(\mathcal{H}\)（它包含 \(\mathcal{C}\)）是什么？
\item
  证明 C 型模的每个短正合列都分裂。描述群 \(\mathrm{K}(\mathcal{C})\)。
\item
  构造一个正合列
\end{enumerate}

\[
0 \to \mathrm{M} _ {0} \to \mathrm{M} _ {1} \to \mathrm{M} _ {2} \to \mathrm{M} _ {3} \to \mathrm{M} _ {4} \to 0
\]

（由 C 型模组成），使得 \(M_{0} \oplus M_{2} \oplus M_{4}\) 与 \(M_{1} \oplus M_{3}\) 在 \(\mathrm{K}(\mathcal{C})\) 中的类不相同。

\begin{enumerate}
\def\labelenumi{\alph{enumi})}
\setcounter{enumi}{3}
\tightlist
\item
  从 \(\mathrm{K}(\mathcal{C})\) 到 \(\mathrm{K}(\mathcal{H})\) 的自然同态是否是单射？
\end{enumerate}

\begin{enumerate}
\def\labelenumi{\arabic{enumi})}
\setcounter{enumi}{1}
\item
  设 \(\mathrm{A}\) 是一个环，\(a\) 与 \(b\) 是 \(\mathrm{A}\) 中满足 \(ab = 1\) 与 \(ba \neq 1\) 的两个元素。证明 A-模 \(\mathrm{A}(1 - ba)\) 是投射的，并且它与零 A-模稳定同构。更精确地说，证明 A-模 \(\mathrm{A}_s\) 与 \(\mathrm{A}(1 - ba) \oplus \mathrm{A}_s\) 同构。
\item
  设 A 与 B 是环。用 A 与 B 的格罗滕迪克群确定格罗滕迪克群 \(R(A \times B)\) 与 \(K_{0}(A \times B)\)（都是 \(A \times B\) 的）。
\item
  设 A 与 B 是森田等价的环。证明群 R(A) 与 R(B) 同构。对 \(K_{0}(A)\) 与 \(K_{0}(B)\) 同样处理。
\item
  设 A 是一个环。用 \(\mathcal{I}(A)\) 表示由偶 \((n,p)\)（其中 n 是整数，p 是 \(M_{n}(A)\) 中的一个幂等元）组成的集。在 \(\mathcal{I}(A)\) 上定义一个加法：令 \((m,p)+(n,q)=(m+n,\begin{pmatrix} p & 0 \\ 0 & q \end{pmatrix})\)。
\end{enumerate}

\begin{enumerate}
\def\labelenumi{\alph{enumi})}
\item
  设 \(\mathcal{P}(A^{\circ})\) 是有限生成投射右 A-模的类所成的幺半群。设 \(\varphi: \mathcal{I}(A) \to \mathcal{P}(A^{\circ})\) 是由 \(\varphi(n,p) = \operatorname{cl}(\operatorname{Im}p)\) 定义的映射。证明 \(\varphi\) 是一个满同态，且元素 \((m,p)\) 与 \((n,q)\)（都属于 \(\mathcal{I}(A)\)）在 \(\varphi\) 下有相同的像当且仅当存在 \(a \in M_{m,n}(A)\) 与 \(b \in M_{n,m}(A)\)，使得 ab = p 且 ba = q。
\item
  由此推出群 \(K_{0}(A)\) 与 \(K_{0}(A^{\circ})\) 同构。
\end{enumerate}

\begin{enumerate}
\def\labelenumi{\arabic{enumi})}
\setcounter{enumi}{5}
\tightlist
\item
  设 A 是一个环，G 是一个交换群，T: A → G 是一个迹型映射（VIII, 第 115 页，习题 3）。对每个有限生成投射右 A-模 P，用
\end{enumerate}

\[
\mathrm{T} _ {\mathrm{P}}: \operatorname{End} _ {\mathrm{A}} (\mathrm{P} \oplus \mathrm{A} _ {d}) \longrightarrow \mathrm{G}
\]

表示扩张 T 的唯一迹型映射（同前所引），并用 \(\pi_{P}\) 表示 \(\operatorname{End}_{\mathrm{A}}(\mathrm{P} \oplus \mathrm{A}_{d})\) 中以 P 为像、以 \(A_{d}\) 为核的投影子。构造一个群同态 \(\widetilde{\mathrm{T}} : \mathrm{K}_{0}(\mathrm{A}) \to \mathrm{G}\)，使得对每个有限生成投射右 A-模 P，有 \(\widetilde{\mathrm{T}}([\mathrm{P}]) = \mathrm{T}_{\mathrm{P}}(\pi_{\mathrm{P}})\)。

\begin{enumerate}
\def\labelenumi{\arabic{enumi})}
\setcounter{enumi}{6}
\tightlist
\item
  设 I 是一个右有向集，\((\mathrm{A}_{i}, f_{ji})\) 是环的一个直接系统，A 是它的极限。证明 \(\mathrm{K}_{0}(\mathrm{A})\) 是群的直接系统 \((\mathrm{K}_{0}(\mathrm{A}_{i}), f_{ji}^{*})\) 的极限。
\end{enumerate}

¶ 8) 设 A、B 与 C 是环，\(f: \mathrm{A} \to \mathrm{C}\) 与 \(g: \mathrm{B} \to \mathrm{C}\) 是环同态；假设 \(f\) 是满射。考虑 \(\mathrm{A} \times \mathrm{B}\) 中由偶 \((a, b)\)（满足 \(f(a) = g(b)\)）组成的子环 D。设 M 是一个 A-模，N 是一个 B-模，\(u\) 是从 \(\mathrm{C} \otimes_{\mathrm{A}} \mathrm{M}\) 到 \(\mathrm{C} \otimes_{\mathrm{B}} \mathrm{N}\) 的一个 C-同构。用 P 表示 \(\mathrm{M} \times \mathrm{N}\) 中由偶 \((m, n)\)（满足 \(u(1 \otimes m) = 1 \otimes n\)）组成的 D-子模。

\begin{enumerate}
\def\labelenumi{\alph{enumi})}
\item
  假设模 M 与 N 分别具有有限基 \((e_{i})_{i\in\mathrm{I}}\) 与 \((f_{j})_{j\in\mathrm{J}}\)，并且 u 在基 \((1\otimes e_{i})\) 与 \((1\otimes f_{j})\) 下的矩阵是一个以 A 为元素的逆矩阵 \((a_{ji})\) 的像。证明 D-模 P 是自由的；更精确地说，元素 \((e_{i},\sum_{j}a_{ji}f_{j})_{i\in\mathrm{I}}\) 构成 P 的一个基。
\item
  设 p 与 q 是整数，\(X \in \mathbf{M}_{p,q}(\mathrm{C})\) 与 \(Y \in \mathbf{M}_{q,p}(\mathrm{C})\) 是满足 \(XY = I_{p}\) 与 \(YX = I_{q}\) 的矩阵。证明方阵 \(\begin{pmatrix} X & 0 \\ 0 & Y \end{pmatrix}\) 是某个以 A 为元素的逆矩阵的像，利用恒等式
\end{enumerate}

\[
\left( \begin{array}{c c} X & 0 \\ 0 & Y \end{array} \right) = \left( \begin{array}{c c} I _ {p} & X \\ 0 & I _ {q} \end{array} \right) \left( \begin{array}{c c} I _ {p} & 0 \\ - Y & I _ {q} \end{array} \right) \left( \begin{array}{c c} I _ {p} & X \\ 0 & I _ {q} \end{array} \right) \left( \begin{array}{c c} 0 & - I _ {p} \\ I _ {q} & 0 \end{array} \right).
\]

\begin{enumerate}
\def\labelenumi{\alph{enumi})}
\setcounter{enumi}{2}
\item
  假设模 M 与 N 是自由且有限生成的。由 a) 与 b) 推出 D-模 P 是投射且有限生成的。
\item
  假设 M 与 N 是投射且有限生成的。构造一个 A-模 \(M'\) 与一个 B-模 \(N'\)，使得 \(M \oplus M'\) 与 \(N \oplus N'\) 是自由且有限生成的，并且 C-模 \(C \otimes_{A} M'\) 与 \(C \otimes_{B} N'\) 同构。由此推出 D-模 P 是投射且有限生成的。
\item
  证明由两个投射导出的典范同态 \(A \otimes_{D} P \to M\) 与 \(B \otimes_{D} P \to N\) 是双射（如上那样化归为 a) 的条件成立的情形）。
\item
  反之，设 \(P'\) 是一个有限生成投射 D-模；令 \(M' = A \otimes_{D} P'\) 与 \(N' = B \otimes_{D} P'\)，并用 \(u' : C \otimes_{A} M' \to C \otimes_{B} N'\) 表示典范同构。证明从 \(P'\) 到 \(M' \times N'\) 的典范同态诱导一个从 \(P'\) 到与 \((M', N', u')\) 相伴的 D-模上的 D-线性同构。
\end{enumerate}

\(g)\) 用 \(\varphi_0:\mathrm{K}_0(\mathrm{D})\to \mathrm{K}_0(\mathrm{A})\times \mathrm{K}_0(\mathrm{B})\) 与 \(\psi_0:\mathrm{K}_0(\mathrm{A})\times \mathrm{K}_0(\mathrm{B})\to \mathrm{K}_0(\mathrm{C})\) 表示满足

\[
\varphi_ {0} ([ \mathrm{P} ]) = ([ \mathrm{A} \otimes_ {\mathrm{D}} \mathrm{P} ], [ \mathrm{B} \otimes_ {\mathrm{D}} \mathrm{P} ]) \quad \text { and } \quad \psi_ {0} ([ \mathrm{M} ], [ \mathrm{N} ]) = [ \mathrm{C} \otimes_ {\mathrm{A}} \mathrm{M} ] - [ \mathrm{C} \otimes_ {\mathrm{B}} \mathrm{N} ].
\]

证明序列

\[
\mathrm{K} _ {0} (\mathrm{D}) \xrightarrow {\varphi_ {0}} \mathrm{K} _ {0} (\mathrm{A}) \times \mathrm{K} _ {0} (\mathrm{B}) \xrightarrow {\psi_ {0}} \mathrm{K} _ {0} (\mathrm{C})
\]

是正合的。证明若此外还存在一个环同态 \(s: C \to A\) 使得 \(f \circ s = Id_{C}\)，则 \(\varphi_{0}\) 是单射且 \(\psi_{0}\) 是满射。

\begin{enumerate}
\def\labelenumi{\arabic{enumi})}
\setcounter{enumi}{8}
\tightlist
\item
  设 A 是一个伪环，即一个结合的但不必含单位元的 Z-代数，\(\widetilde{A}\) 是由 A 通过添加单位元（III, §1, No.~2, 第 431 页）得到的 Z-代数；集合 \(\widetilde{\mathrm{A}}\) 等于 \(\mathbf{Z} \times \mathrm{A}\)。把同态 \((n, a) \mapsto n\) 记作 \(\varepsilon : \widetilde{\mathrm{A}} \to \mathbf{Z}\)，并把同态 \(\varepsilon^{*} : \mathrm{K}_{0}(\widetilde{\mathrm{A}}) \to \mathrm{K}_{0}(\mathbf{Z})\)（它与 \(\varepsilon\) 相伴）的核记作 \(\mathrm{K}_{0}(\mathrm{A})\)。
\end{enumerate}

\begin{enumerate}
\def\labelenumi{\alph{enumi})}
\item
  证明若 A 含单位元，则 \(\tilde{A}\) 同构于积环 \(Z \times A\)，且 \(K_{0}(A)\) 的这一定义与 No.~8 中给出的定义相容。
\item
  设 A 与 B 是伪环，\(f: \mathrm{A} \to \mathrm{B}\) 是一个同态。设 \(\widetilde{f}: \widetilde{\mathrm{A}} \to \widetilde{\mathrm{B}}\) 是由 \(\widetilde{f}((n, a)) = (n, f(a))\) 定义的环同态。证明 \(\widetilde{f}^*\) 把 \(\mathrm{K}_0(\mathrm{A})\) 映到 \(\mathrm{K}_0(\mathrm{B})\) 之中。把由 \(\widetilde{f}^*\) 导出的同态记作 \(f^*: \mathrm{K}_0(\mathrm{A}) \to \mathrm{K}_0(\mathrm{B})\)。c) 假设 \(f\) 是满射；用 \(\mathfrak{a}\) 表示它的核，并把 \(\mathfrak{a}\) 到 A 中的典范单射记作 \(i\)。证明序列
\end{enumerate}

\[
\mathrm{K} _ {0} (\mathfrak {a}) \xrightarrow {i ^ {*}} \mathrm{K} _ {0} (\mathrm{A}) \xrightarrow {f ^ {*}} \mathrm{K} _ {0} (\mathrm{B})
\]

是正合的。若此外还存在伪环的同态 \(s: B \to A\) 使得 \(f \circ s = Id_{B}\)，则 \(i^{*}\) 是单射且 \(f^{*}\) 是满射（对三元组 (A, B, C) 先取 (A, A, B) 再取 (D, Z, A)，应用习题 8）。

\begin{enumerate}
\def\labelenumi{\alph{enumi})}
\setcounter{enumi}{3}
\tightlist
\item
  证明习题 7 与 8 的结论可推广到伪环的情形。e) 用 \(\mathbf{M}_{\infty}(\mathrm{A})\) 表示以 A 为元素且仅有有限多个非零元素的 \(\mathbf{N} \times \mathbf{N}\) 矩阵所成的伪环。设 \(f: \mathrm{A} \to \mathbf{M}_{\infty}(\mathrm{A})\) 是把 \(a \in \mathrm{A}\) 映到矩阵 \((b_{ij})\)（满足 \(b_{00} = a\) 与 \(b_{ij} = 0\)，当 \((i,j) \neq (0,0)\)）的映射。证明 \(f^{*}\) 是从 \(\mathrm{K}_0(\mathrm{A})\) 到 \(\mathrm{K}_0(\mathbf{M}_{\infty}(\mathrm{A}))\) 的同构（利用习题 4 与 7）。
\end{enumerate}

\begin{enumerate}
\def\labelenumi{\arabic{enumi})}
\setcounter{enumi}{9}
\tightlist
\item
  设 A 是一个环。用 \(\mathbf{GL}_{\infty}(\mathrm{A})\) 表示 \(\mathbf{GL}(\mathrm{A}^{(\mathbf{N})})\) 中由满足 X - I 仅有有限多个非零元素的矩阵 X 组成的子群，并用 \(\mathrm{E}(\mathrm{A})\) 表示 \(\mathbf{GL}_{\infty}(\mathrm{A})\) 中由形如 \(I + \lambda E_{ij}\)（其中 i, j 属于 N，\(i \neq j\)，\(\lambda \in A\)）的矩阵生成的子群；它是 \(\mathbf{GL}_{\infty}(\mathrm{A})\) 的导群（II, §10, 第 422 页，习题 15）。把商群 \(\mathbf{GL}_{\infty}(\mathrm{A}) / \mathrm{E}(\mathrm{A})\) 记作 \(\mathrm{K}_{1}(\mathrm{A})\)；它的群运算写作加法。每个环同态 \(f : A \to B\) 诱导一个群同态 \(f^{*} : K_{1}(A) \to K_{1}(B)\)。
\end{enumerate}

\begin{enumerate}
\def\labelenumi{\alph{enumi})}
\item
  对每个整数 \(n \geqslant 1\)，设 \(\iota_{n} : \mathbf{GL}_{n}(\mathrm{A}) \to \mathrm{K}_{1}(\mathrm{A})\) 是把 \(X \in \mathbf{GL}_{n}(\mathrm{A})\) 映到矩阵 \(\begin{pmatrix} X & 0 \\ 0 & I \end{pmatrix}\) 的类的同态。证明若环 A 是局部环，则对每个 n，\(\iota_{n}\) 是满射（利用 VIII, 第 42 页，习题 17）。
\item
  假设环 A 是交换的。构造一个满同态 \(d_{\mathrm{A}} : \mathrm{K}_{1}(\mathrm{A}) \to \mathrm{A}^{*}\)，使得对每个 n 有 \(d_{A} \circ \iota_{n} = det\)。若此外 A 是局部环或欧几里得环，则同态 \(d_{A}\) 与 \(\iota_{1}\) 是互逆的双射。
\item
  设 I 是一个右有向集，\((\mathrm{A}_{i}, f_{ji})\) 是环的一个直接系统，A 是它的极限。证明 \(\mathrm{K}_{1}(\mathrm{A})\) 是群的直接系统 \((\mathrm{K}_{1}(\mathrm{A}_{i}), f_{ji}^{*})\) 的极限。
\end{enumerate}

¶ 11) 采用习题 8 的假设与记号。用 \(f': \mathrm{D} \to \mathrm{B}\) 与 \(g': \mathrm{D} \to \mathrm{A}\) 表示由两个投射导出的同态，并用

\[
\varphi_ {1}: \mathrm{K} _ {1} (\mathrm{D}) \rightarrow \mathrm{K} _ {1} (\mathrm{A}) \times \mathrm{K} _ {1} (\mathrm{B}) \quad \text { and } \quad \psi_ {1}: \mathrm{K} _ {1} (\mathrm{A}) \times \mathrm{K} _ {1} (\mathrm{B}) \rightarrow \mathrm{K} _ {1} (\mathrm{C})
\]

表示由 \(\varphi_1(z) = (g'^* (z), f'^* (z))\) 与 \(\psi_1(x,y) = f^* (x) - g^* (y)\) 定义的映射。

\begin{enumerate}
\def\labelenumi{\alph{enumi})}
\tightlist
\item
  证明序列
\end{enumerate}

\[
\mathrm{K} _ {1} (\mathrm{D}) \xrightarrow {\varphi_ {1}} \mathrm{K} _ {1} (\mathrm{A}) \times \mathrm{K} _ {1} (\mathrm{B}) \xrightarrow {\psi_ {1}} \mathrm{K} _ {1} (\mathrm{C})
\]

是正合的（注意 \(f(\mathrm{E}(\mathrm{A})) = \mathrm{E}(\mathrm{C})\)）。证明若此外还存在一个环同态 \(s : C \to A\) 使得 \(f \circ s = Id_{C}\)，则 \(\psi_{1}\) 是满射。

\begin{enumerate}
\def\labelenumi{\alph{enumi})}
\setcounter{enumi}{1}
\item
  对每个整数 n 与 \(\mathbf{GL}_{n}(\mathrm{C})\) 的每个元素 g，用 \(P_{g}\) 表示与模 \(A^{n}\)、\(B^{n}\) 以及从 \(C \otimes_{A} A^{n}\) 到 \(C \otimes_{B} B^{n}\) 的以 g 为矩阵的同构相伴的投射 D-模（习题 8）。构造一个同态 \(\partial : K_{1}(\mathrm{C}) \to K_{0}(\mathrm{D})\)，使得 \(\partial \circ \iota_{n}(g) = [\mathrm{P}_{g}] - [\mathrm{D}_{s}^{n}]\) 对每个整数 n 与 \(\mathbf{GL}_{n}(\mathrm{C})\) 的每个元素 g 成立。
\item
  证明序列
\end{enumerate}

\[
\mathrm{K} _ {1} (\mathrm{D}) \xrightarrow {\varphi_ {1}} \mathrm{K} _ {1} (\mathrm{A}) \times \mathrm{K} _ {1} (\mathrm{B}) \xrightarrow {\psi_ {1}} \mathrm{K} _ {1} (\mathrm{C}) \xrightarrow {\partial} \mathrm{K} _ {0} (\mathrm{D}) \xrightarrow {\varphi_ {0}} \mathrm{K} _ {0} (\mathrm{A}) \times \mathrm{K} _ {0} (\mathrm{B}) \xrightarrow {\psi_ {0}} \mathrm{K} _ {0} (\mathrm{C})
\]

是正合的。

\begin{enumerate}
\def\labelenumi{\arabic{enumi})}
\setcounter{enumi}{11}
\tightlist
\item
  设 \(f: \mathrm{A} \to \mathrm{B}\) 是一个满射环同态，\(\mathfrak{a}\) 是其核，且 \(i\) 是从 \(\mathfrak{a}\) 到 \(\mathrm{A}\) 的典范单射。构造一个同态 \(\partial: \mathrm{K}_1(\mathrm{B}) \to \mathrm{K}_0(\mathfrak{a})\)，使得序列
\end{enumerate}

\[
\mathrm{K} _ {1} (\mathrm{A}) \xrightarrow {f ^ {*}} \mathrm{K} _ {1} (\mathrm{B}) \xrightarrow {\partial} \mathrm{K} _ {0} (\mathfrak {a}) \xrightarrow {i ^ {*}} \mathrm{K} _ {0} (\mathrm{A}) \xrightarrow {f ^ {*}} \mathrm{K} _ {0} (\mathrm{B})
\]

是正合的（按习题 11, c) 的方式论证）。

\begin{enumerate}
\def\labelenumi{\arabic{enumi})}
\setcounter{enumi}{12}
\item
  用 \(\mathcal{C}\)（相应地，\(\mathcal{D}\)）表示有限长度 \(\mathbf{Z}\)-模（相应地，有限生成 \(\mathbf{Z}\)-模）的类所成之集。证明群同态 \(\gamma_{\mathcal{D},\mathcal{C}}\)（从 \(\mathrm{K}(\mathcal{C})\) 到 \(\mathrm{K}(\mathcal{D})\)）是零同态。
\item
  设 K 是一个特征异于 2 的交换域。令
\end{enumerate}

\[
\mathrm{A} = \mathrm{K} [ \mathrm{X}, \mathrm{Y}, \mathrm{Z} ] / (\mathrm{X} ^ {2} + \mathrm{Y} ^ {2} + \mathrm{Z} ^ {2} - 1).
\]

设 M 是 A-模 \(\mathrm{A}^{3}/\mathrm{A}(\overline{\mathrm{X}},\overline{\mathrm{Y}},\overline{\mathrm{Z}})\)，其中 \(\overline{X}\)、\(\overline{Y}\) 与 \(\overline{Z}\) 分别表示 X、Y 与 Z 在 A 中的像。证明 \(M\oplus A\) 是自由 A-模，但 M 不是自由模。

